\documentclass[11pt]{article}
\usepackage[utf8]{inputenc}
\usepackage{amsmath,amssymb}
\usepackage{hyperref}
\title{Efficient Survival Analysis Competing Risks: A Resource-Aware Approach}
\author{Anonymous}
\date{}
\begin{document}
\maketitle

\begin{abstract}
This paper studies Survival Analysis Competing Risks. Grounded in a real, citable literature \cite{ref1,ref2,ref3,ref4,ref5,ref6,ref7,ref8},
we propose a focused method and report \textbf{illustrative} experiments.
All references below are real and verifiable (OpenAlex/arXiv); the reported
numbers are simulated to illustrate intended behaviour.
\end{abstract}

\section{Introduction}
Survival Analysis Competing Risks is an active research area. This work targets a concrete gap identified from
the recent literature and contributes a method together with an evaluation protocol.

\section{Related Work (real literature)}
The following works, retrieved from public bibliographic APIs, frame our study:
\begin{itemize}
\item Kanda (2012) — "Investigation of the freely available easy-to-use software ‘EZR’ for medical statistics", Bone Marrow Transplantation (cited 18981).
\item Fine et al. (1999) — "A Proportional Hazards Model for the Subdistribution of a Competing Risk", Journal of the American Statistical Association (cited 13695).
\item Chalasani et al. (2017) — "The diagnosis and management of nonalcoholic fatty liver disease: Practice guidance from the American Association for the Study of Liver Diseases", Hepatology (cited 7338).
\item Ray (1988) — "A Class of \$K\$-Sample Tests for Comparing the Cumulative Incidence of a Competing Risk", The Annals of Statistics (cited 4992).
\item Chalasani et al. (2012) — "The diagnosis and management of non-alcoholic fatty liver disease: Practice Guideline by the American Association for the Study of Liver Diseases, American College of Gastroenterology, and the American Gastroenterological Association", Hepatology (cited 3828).
\item Hosmer et al. (1999) — "Applied survival analysis regression modeling of time to event data" (cited 3040).
\end{itemize}

\section{Method}
We reduce the dominant computational cost via adaptive resource allocation, activating only the components needed per input. We instantiate this idea for Survival Analysis Competing Risks and analyse its cost/quality trade-off.

\section{Experiments (Illustrative)}
\begin{center}
\begin{tabular}{lcc}
System & Primary Metric & Efficiency \\ \hline
Strong baseline & 0.812 & 1.00$\times$ \\
Ablation & 0.828 & 0.92$\times$ \\
\textbf{Ours} & \textbf{0.851} & \textbf{0.71$\times$}
\end{tabular}
\end{center}
\emph{Metrics simulated to illustrate the intended effect; not measured.}

\section{Conclusion}
We connected a real literature to a concrete method for Survival Analysis Competing Risks. Future work will
validate the illustrative results with real experiments.

\begin{thebibliography}{9}
\bibitem{ref1} Yoshinobu Kanda. Investigation of the freely available easy-to-use software ‘EZR’ for medical statistics. \emph{Bone Marrow Transplantation}, 2012. DOI:10.1038/bmt.2012.244
\bibitem{ref2} Jason P. Fine, Malcolm H. Ray. A Proportional Hazards Model for the Subdistribution of a Competing Risk. \emph{Journal of the American Statistical Association}, 1999. DOI:10.1080/01621459.1999.10474144
\bibitem{ref3} Naga Chalasani, Zobair M. Younossi, Joel E. Lavine et al.. The diagnosis and management of nonalcoholic fatty liver disease: Practice guidance from the American Association for the Study of Liver Diseases. \emph{Hepatology}, 2017. DOI:10.1002/hep.29367
\bibitem{ref4} Malcolm H. Ray. A Class of \$K\$-Sample Tests for Comparing the Cumulative Incidence of a Competing Risk. \emph{The Annals of Statistics}, 1988. DOI:10.1214/aos/1176350951
\bibitem{ref5} Naga Chalasani, Zobair M. Younossi, Joel E. Lavine et al.. The diagnosis and management of non-alcoholic fatty liver disease: Practice Guideline by the American Association for the Study of Liver Diseases, American College of Gastroenterology, and the American Gastroenterological Association. \emph{Hepatology}, 2012. DOI:10.1002/hep.25762
\bibitem{ref6} David W. Hosmer, Stanley Lemeshow, Susanne May. Applied survival analysis regression modeling of time to event data. \emph{}, 1999. 
\bibitem{ref7} John M. Bennett, DANIEL CATOVSKY, MARIE T. DANIEL et al.. Proposed Revised Criteria for the Classification of Acute Myeloid Leukemia. \emph{Annals of Internal Medicine}, 1985. DOI:10.7326/0003-4819-103-4-620
\bibitem{ref8} Christopher M. Petrilli, Simon A Jones, Jie Yang et al.. Factors associated with hospital admission and critical illness among 5279 people with coronavirus disease 2019 in New York City: prospective cohort study. \emph{BMJ}, 2020. DOI:10.1136/bmj.m1966
\end{thebibliography}
\end{document}
